\subsection{The entire global zero-first-coordinate family in the inherited module quotient}

\begin{statement}{E.47}{Proved}{Global section coordinates from explicit norm-coherent cyclotomic units}
Use the actual arithmetic group $H$ and the unit tangential
origin of (E38.5)--(E38.8). Write
\[
 \mathscr A=\mathbb Z_3[[\mathbb Z_3]]
       =\mathbb Z_3[[U]],\qquad T=1+U,\qquad
 J=(3,U)\mathscr A,\qquad
 H=3\mathbb Z_3\ltimes J.
 \tag{E47.1}
\]
Here $T^x=\delta_x$, multiplication by $T$ translates
the measure variable, and
$g(a,v)=(\chi(g)a,\chi(g)v(T^{\chi(g)}))$.
Let $\operatorname{Sec}_0(H)$ denote the geometric
conjugacy classes of global sections whose image in
$H^1(\mathbb Q,3\mathbb Z_3(1))$ is zero. There is an
explicit isomorphism of pointed topological spaces,
and a group structure specified by its right side,
\[
 \operatorname{Sec}_0(H)
 \xrightarrow{\sim}
 3\widehat{\mathbb Q^{\times}}_3
       \ \times\prod_{j\ge0}
       \left(\mathbb Z_3(1)\oplus
                    \mathbb Z_3[[\Gamma]]\right),
 \qquad \Gamma=1+3\mathbb Z_3.
 \tag{E47.2}
\]
The completion means exactly
$\widehat{F^\times}_3=\varprojlim_mF^\times/(F^\times)^{3^m}$.
In particular the first factor is the group of restricted
prime-valuation sequences
\[
 3\widehat{\mathbb Q^\times}_3
  =\left\{(3b_p)_p\in\prod_{p\ {\rm prime}}3\mathbb Z_3:
       \text{for every }m,\ b_p\in3^m\mathbb Z_3
                 \text{ for all but finitely many }p\right\},
 \tag{E47.3}
\]
with its inverse-limit topology, not an asserted compact
topology. The factors indexed by $j$ in (E47.2) are
compact; their displayed product is the full product.

The coordinate maps are as follows. First conjugate a
section so its first cocycle is identically zero. This
normalization is unique in the first geometric coordinate.
Its second coordinate is a cocycle $v:G_{\mathbb Q}\to J$.
Take its augmentation $\nu(g)=v(g)(1)$, and for each $j\ge0$
restrict its measure to $3^j\mathbb Z_3^\times$ and scale
that set by $3^{-j}$. Under the exact Shapiro and norm
maps of E.46, the latter class is the compatible system
\[
 (u_{j,n})_{n\ge1}\in
       \varprojlim_{N_{K_{n+1}/K_n}}\widehat{K_n^\times}_3,
 \qquad K_n=\mathbb Q(\zeta_n),\quad
 \zeta_{n+1}^3=\zeta_n,\quad \operatorname{ord}(\zeta_n)=3^n.
 \tag{E47.4}
\]
Every such system has a unique expression
\[
 \begin{split}
 u_{j,n}&=\zeta_n^{\,a_j\bmod 3^n}
        \prod_{\gamma\in\Gamma_n}
                \gamma(\epsilon_n)^{\,b_{j,n,\gamma}},\\
 \epsilon_n&=(1-\zeta_n)(1-\zeta_n^{-1}),\qquad
 \Gamma_n=\Gamma/(1+3^n\mathbb Z_3),\\
 a_j&\in\mathbb Z_3,\qquad
 b_j=(\sum_{\gamma\in\Gamma_n}
             b_{j,n,\gamma}[\gamma])_n
                   \in\mathbb Z_3[[\Gamma]].
 \end{split}
 \tag{E47.5}
\]
Three-adic powers are taken in the indicated
multiplication-by-$3^m$ completion. Equations
(E47.4)--(E47.5), together with $\nu$, are the coordinates
in (E47.2). The root-of-unity line has its usual Tate
action under the natural field action of $\Gamma$.
On measures this action is the inverse deck action
$\delta_x\mapsto\delta_{\gamma^{-1}x}$; the proof
specifies the resulting involution on the group ring.

The two-coordinate augmentation map
$H\to(3\mathbb Z_3(1))^2$,
$(a,v)\mapsto(a,\operatorname{aug}v)$, and the
scaled toric map (E46.7) have, respectively, the
following exact maps on the displayed parameters:
\[
 \begin{aligned}
 (\nu,(a_j,b_j)_{j\ge0})
       &\longmapsto (0,\nu)
          \in H^1(\mathbb Q,(3\mathbb Z_3(1))^2),\\
 (\nu,(a_j,b_j)_{j\ge0})
       &\longmapsto (0,\nu/3)
          \in H^1(\mathbb Q,\mathbb Z_3(1)^2).
 \end{aligned}
 \tag{E47.6}
\]
Consequently every such toric coordinate occurs, and
each fiber is the displayed compact product in (E47.2).
Every annular coordinate is unramified outside $3$.
Its independent arithmetic evaluation is the explicit
free-group presentation
\[
 \begin{split}
 H^1(\mathbb Q,
       \mathbb Z_3(1)\otimes\mathbb Z_3[[\mathbb Z_3^\times]])
 &\simeq
 \frac{B\oplus B}{\{(d z,d z):z\in B\}}
       \xrightarrow{\sim} B\oplus\mathbb Z_3,\\
 B&=\mathbb Z_3[[\Gamma]],\qquad d=7S_7-1,\\
 [(a,b)]&\longmapsto
       \bigl(a-b,\ b\big|_{S_7=1/7}\bigr).
 \end{split}
 \tag{E47.7}
\]
In this formula $S_\gamma$ denotes direct positional
translation $\delta_x\mapsto\delta_{\gamma x}$.
Thus the natural field action on the last line is
$S_\gamma^{-1}$, acting as multiplication by $\gamma$.

The actual geometric maps being evaluated are
$N\twoheadrightarrow H$ with kernel
$J_y=[I_y,I_y]$, and
\[
 N/N''\twoheadrightarrow H,\quad\ker=J_y/N'',\qquad
 N/\gamma_3I_{\rm marked}\twoheadrightarrow H,\quad
                 \ker=J_y/\gamma_3I_{\rm marked},
 \tag{E47.8}
\]
for the markings in (E38.9). The parameterization
(E47.2) evaluates the target of these precise maps.

\emph{Manuscript: \S51.16, folios733v--736v,
scans172--178, (371)--(379), Part E(c),(d).
The finite program \texttt{e\_global\_zero.py}
checks the cyclotomic elements, their actual norm
maps, and the compatible finite presentations
underlying (E47.7).}
\end{statement}
\begin{proof}
\emph{Normalization, the centered decomposition, and topology.}
If the first cocycle has zero class, it equals
$g\mapsto(\chi(g)-1)s$ for some $s\in3\mathbb Z_3$.
Conjugation by $(s,0)$ makes it zero. This $s$ is
unique: complex conjugation acts by $-1$, and
$3\mathbb Z_3$ has no two-torsion. In fact
$s=-a(c)/2$, so this normalization is continuous.
Once the first
cocycle is zero, a geometric conjugation preserving
it has first coordinate zero and changes $v$ by the
ordinary coboundary of an element of $J$. Thus
the normalized section set is exactly the abelian
cohomology quotient $H^1(\mathbb Q,J)$, including its
uniform continuous-cochain quotient topology.

There is an actual continuous $G_{\mathbb Q}$-module
isomorphism
\[
 \begin{split}
 J&\longrightarrow 3\mathbb Z_3(1)\delta_0
       \oplus\prod_{j\ge0}M_j,\\
 \mu&\longmapsto
       (\operatorname{aug}\mu,\,
                    (\mu|_{3^j\mathbb Z_3^\times})_j),\\
 M_j&=\mathbb Z_3(1)\otimes
                    \mathbb Z_3[[3^j\mathbb Z_3^\times]].
 \end{split}
 \tag{E47.9}
\]
Its inverse is the centered series
\[
 (\nu,(\mu_j)_j)\longmapsto
      \nu\delta_0+
           \sum_{j\ge0}(\mu_j-\operatorname{aug}(\mu_j)\delta_0).
 \tag{E47.10}
\]
In the positional quotient $\mathbb Z/3^k$ every
summand with $j\ge k$ is zero. Hence this series
converges in the completed group-ring topology for
arbitrary choices of its factors; no sum of their
uncentered augmentations is required. Restricting
the sum to each annulus recovers that factor, and
its augmentation is $\nu$. Conversely a measure of
augmentation zero whose restrictions to all the
annuli vanish has support at zero and is zero.
These checks prove both inverse identities and
continuity. Multiplication of the variable by a
cyclotomic unit preserves every annulus, so the
maps are equivariant, including their Tate twist.
Scaling identifies each $M_j$ with the module $M$
in (E47.7).

Continuous cochains into a product are precisely
products of continuous cochains, and the uniform
cochain topology into that product is the product
of the uniform topologies. Cocycle equations and
coboundaries are componentwise; the quotient maps
are open. Therefore (E47.9) induces an isomorphism
of topological cohomology groups with the product
of the cohomology groups of its factors. The first
factor is $3\widehat{\mathbb Q^\times}_3$ by Kummer
theory. Indeed $\mathbb Q^\times/\{\pm1\}$ is the
free abelian group on the rational primes, and
completion gives precisely (E47.3). The inclusion
$3\mathbb Z_3(1)\hookrightarrow\mathbb Z_3(1)$
induces multiplication by three on this group.

\emph{The exact annular Shapiro maps.}
At positional level $n\ge1$, the unit positions
are $(\mathbb Z/3^n)^\times$. Their stabilizer at
one is $G_{K_n}$. Shapiro evaluation at that
position gives
\[
 H^1(\mathbb Q,
      \mathbb Z_3[(\mathbb Z/3^n)^\times](1))
       \xrightarrow{\sim}\widehat{K_n^\times}_3.
 \tag{E47.11}
\]
Projection of unit positions induces corestriction;
the Kummer formula for corestriction is the product
of conjugate roots, so the map is exactly the field
norm. E.46 proves the inverse-limit isomorphism
by aligning continuous cocycles and coboundaries.
The relevant degree-zero groups vanish because
no number field contains a compatible sequence
of roots of unity of unbounded three-power order.
That proof does not require compactness of
$\widehat{K_n^\times}_3$.

Here is the action convention explicitly. For a
cocycle $c$ at positional level $n$, a lift $g$
with $\chi(g)=\gamma$, and $h\in G_{K_n}$, the
cocycle identity gives
\[
 \gamma\,c(g^{-1}hg)_1
       =c(h)_\gamma+(\chi(h)-1)c(g)_\gamma.
 \tag{E47.12}
\]
The last term is a coboundary. Consequently the
natural field action by $\gamma$ corresponds to
the inverse deck translation, whose coefficient
at one is $c(h)_\gamma$. The comparison between
the direct positional ring and the group ring
in (E47.5) is the explicit involution
$S_\gamma\mapsto[\gamma^{-1}]$.

\emph{The three-primary class groups in this tower vanish.}
The ring $\mathbb Z[\zeta_1]$ is Euclidean for
its norm, so $K_1$ has class number one.
For $K_{n+1}/K_n$ there is one ramified finite
prime, the unique prime above three, and it is
totally ramified. There is no infinite-place
ramification. These facts follow directly from
$\Phi_{3^n}(X)=X^{2\cdot3^{n-1}}+X^{3^{n-1}}+1$:
$\Phi_{3^n}(X+1)$ is Eisenstein at three, and
the polynomial discriminant is, up to sign,
$3^{(2n-1)3^{n-1}}$. The latter identity follows
by taking norms of
$\Phi_{3^n}'(\zeta_n)=
3^n\zeta_n^{-1}/(\zeta_n^{3^{n-1}}-1)$.
The field discriminant divides it, excluding
ramification at every other finite prime.

Apply the exact ambiguous class number formula
of \cite[Theorem 1, p.~1 of the opened preprint]{E_LemmermeyerAmbiguous}
to this cyclic extension of degree three. Its
ramification factor is $3^{1-1}=1$, so
\[
 |\operatorname{Cl}(K_{n+1})^{C_3}|
   =\frac{h(K_n)}
    {[\mathcal O_{K_n}^{\times}:
       \mathcal O_{K_n}^{\times}\cap
                 N_{K_{n+1}/K_n}K_{n+1}^{\times}]}.
 \tag{E47.13}
\]
If $3\nmid h(K_n)$, the integer on the left
is prime to three. A nonzero finite three-group
acted on by $C_3$ has a nonzero fixed element:
nonfixed orbits have size three, so the number
of fixed elements is divisible by three and
already includes zero. Applied to the three-primary class
group, this proves it is zero. Induction proves
$3\nmid h(K_n)$ for every $n$.
This is also the one-prime class-group argument
proved in \cite[Lemma 1.3.13]{E_SharifiIwasawa}.

Put $F_n=K_n^+$. Extension of ideals followed
by norm multiplies $\operatorname{Cl}(F_n)$
by two. It is therefore injective on its
three-primary subgroup, whose image in
$\operatorname{Cl}(K_n)$ is zero. Thus
$3\nmid h(F_n)$ as well.

\emph{Every norm-coherent class is a three-unit.}
Fix a rational prime $p\ne3$. Its decomposition
subgroup in $\Gamma$ is open: its positive
cyclotomic component has infinite order, since
no positive integral power of $p$ is $1$ or
$-1$ as a rational, hence as a three-adic number.
After a finite level the number of primes over
$p$ is constant and their residue degrees in
successive layers are three. For a compatible
system $u_n$, norm compatibility therefore
makes the valuation at each fixed finite-level
prime divisible by arbitrarily large powers
of three. That valuation is in $\mathbb Z_3$,
so it is zero.

This valuation argument applies to the completed
multiplicative groups themselves. To verify the
kernel, let $\mathcal U_n=\mathcal O_{K_n}[1/3]^\times$.
The image of the away-from-three divisor map
$K_n^\times\to\bigoplus_{\mathfrak p\nmid3}\mathbb Z$
is a free abelian subgroup with cokernel the
finite $S$-class group, of order prime to three.
The sequence with kernel $\mathcal U_n$ splits
as a sequence of abstract abelian groups because
its image is free. Completion is exact on this
split sequence; the inclusion of that image into
the divisor group induces an isomorphism on
three-completions, since its finite index is a
three-adic unit. Hence the kernel of all the
completed valuations is exactly
$\widehat{\mathcal U_n}_3$. We have proved the
actual equality of inverse limits
\[
 \varprojlim_N\widehat{K_n^\times}_3
       =\varprojlim_N\widehat{\mathcal U_n}_3.
 \tag{E47.14}
\]

\emph{A basis of the completed three-units and its norms.}
At level $n$ the real field has degree
$r_n=3^{n-1}$. The element $\epsilon_n$ generates
its unique prime above three: its valuation
there is one, and its absolute norm is three.
Let $\epsilon_{n,a}=(1-\zeta_n^a)(1-\zeta_n^{-a})$
for the $r_n$ unit classes modulo signs.
The normalized real circular units are
\[
 \eta_{n,a}=\zeta_n^{(1-a)/2}
                     \frac{1-\zeta_n^a}{1-\zeta_n},
 \qquad
              \frac{\epsilon_{n,a}}{\epsilon_n}
                         =\eta_{n,a}^2.
 \tag{E47.15}
\]
The half exponent is read modulo $3^n$.
The opened proof of the Hasse index theorem,
\cite[Theorem 4.3.3 and its proof, pp.~108--109]{E_SharifiCyclotomic},
states that these $r_n-1$ units form a basis
of the real circular units modulo $\{\pm1\}$
and that their index in the full real units
is $h(F_n)$. Thus the subgroup generated by
$\{\pm1\}$ and the $\epsilon_{n,a}$ has the
exact index
\[
       2^{r_n-1}h(F_n)
       \quad\text{in }\mathcal O_{F_n}[1/3]^\times.
 \tag{E47.16}
\]
Indeed $\epsilon_n$ accounts for the valuation
at three, while the ratios in (E47.15) account
for the unit kernel. This index is prime to
three. Both the number of generators and the
rank of the real three-units are $r_n$, so
\[
 \mathbb Z_3[\Gamma_n]\longrightarrow
       \widehat{\mathcal O_{F_n}[1/3]^\times}_3,
 \qquad\sum b_\gamma[\gamma]\longmapsto
                 \prod\gamma(\epsilon_n)^{b_\gamma}
 \tag{E47.17}
\]
is an isomorphism of compact $\mathbb Z_3$-modules.

Complex conjugation gives integral plus and
minus projectors $(1\pm c)/2$ on the
three-completion. Its plus part is precisely
the completion in (E47.17): norm followed by
inclusion acts as two on invariants. For
$u\in\mathcal U_n$, the quotient $u/\bar u$
has valuation zero at every finite prime,
because the only allowed prime is fixed by
conjugation, and every conjugate has complex
absolute value one. It is a root of unity.
For completeness this last assertion follows
by applying boundedness of the coefficients
of minimal polynomials to all its powers:
only finitely many such integral polynomials
occur, so two powers coincide. The minus
part of the completion is consequently
$\mu_{3^n}$. We have proved the actual
direct-sum decomposition
\[
 \widehat{\mathcal U_n}_3
       =\mu_{3^n}\oplus
         \left\langle\gamma(\epsilon_n):
                      \gamma\in\Gamma_n\right\rangle_{\mathbb Z_3}.
 \tag{E47.18}
\]

The relative Galois group has exponents
$1,1+3^n,1+2\cdot3^n$. Multiplication gives
\[
 N\zeta_{n+1}=\zeta_n,\qquad
 N(1-\zeta_{n+1})=1-\zeta_n,\qquad
 N\epsilon_{n+1}=\epsilon_n.
 \tag{E47.19}
\]
The middle identity is also proved as
\cite[Lemma 4.3.6]{E_SharifiCyclotomic}.
Norm commutes with the abelian field action,
so on (E47.18) it is reduction
$\mathbb Z/3^{n+1}\to\mathbb Z/3^n$ on the
root exponents and the group-ring projection
$\mathbb Z_3[\Gamma_{n+1}]\to\mathbb Z_3[\Gamma_n]$
on the real unit exponents. There is no extra
factor of three in this latter map: each
conjugate generator maps to its conjugate
generator downstairs. Taking inverse limits
proves (E47.5), its uniqueness, and
\[
       \varprojlim_N\widehat{K_n^\times}_3
            \simeq\mathbb Z_3(1)\oplus\mathbb Z_3[[\Gamma]].
 \tag{E47.20}
\]
This also constructs the inverse section map.
Use the compatible Kummer cocycles and the
inverse Shapiro map of E.46 for each system
(E47.5), scale their measures to their annuli,
and form (E47.10) together with a Kummer
cocycle representing $\nu$. Every finite
positional quotient sees only finitely many
summands, so this is a continuous $J$-valued
cocycle. The section is $g\mapsto((0,v_g),g)$.
Changes of the chosen representatives are
exactly the componentwise coboundaries in
(E47.9), hence geometric conjugations.

\emph{The independent free-inertia calculation and the topology of the limit.}
The annular module $M$ is the finite-index
induced module from $G_K$, $K=\mathbb Q(\zeta_1)$,
of the module $\mathcal B$ whose underlying
ring is $B=\mathbb Z_3[[\Gamma]]$ and whose action is
$g\cdot z=\chi(g)S_{\chi(g)}z$.
The two cosets of $G_K$ correspond to the
positive and negative residue classes of unit
positions modulo three; on the second coset
complex conjugation also contributes its
Tate sign. This specifies the induction map
and proves the continuous Shapiro comparison
with $H^1(K,\mathcal B)$.

For a prime $\mathfrak q\nmid3$ of $K$, the
coefficient action is unramified. The restriction
of a cocycle to inertia is therefore an additive
homomorphism. Wild inertia and prime-to-three
tame inertia have zero image in this pro-three
module. Put $q=N\mathfrak q$ and let $v$ be
its value at a pro-three tame generator.
The relation $F\tau F^{-1}=\tau^q$ gives
\[
        qS_qv=qv,\qquad (S_q-1)v=0.
 \tag{E47.21}
\]
The ring $B$ is the domain $\mathbb Z_3[[S_7-1]]$;
$S_q-1$ is nonzero since $q\ne1$ in
$1+3\mathbb Z_3$. Thus $v=0$. The cocycle
kills all inertia away from three. Equivalently
it kills every factor of the actual free-inertia
kernel (E39.1), so it factors through
$G_{K,\{3\}}$. Its graph has pro-three target,
and it factors further through the free group
$P=\langle\tau,\tau'\rangle$ of E.37.

Both actual Frobenius generators have cyclotomic
value seven. A $B$-valued cocycle is freely
specified by the pair of its values on them,
and the coboundaries are the diagonal pairs
$((7S_7-1)z,(7S_7-1)z)$. Evaluation at
$S_7=1/7$ is continuous because $1/7-1\in3\mathbb Z_3$.
Its kernel is the principal ideal $(7S_7-1)$:
formal division by $S_7-1-(1/7-1)$ converges
integrally. The kernel and inverse of the
last map of (E47.7) now follow directly; a
representative of $(x,y)$ is $(x+y,y)$,
where $y$ is the constant series. This
proves that presentation and also shows its
cohomology quotient is compact Hausdorff.
Finite-index Shapiro is continuous, so the
annular cohomology has this compact topology.
The natural map to (E47.20) is a continuous
bijection onto an inverse limit of compact
three-unit groups. It is therefore a
homeomorphism. This establishes the claimed
topology without treating the larger groups
$\widehat{K_n^\times}_3$ as compact.

Finally, (E47.9) sends augmentation exactly
to its first summand, proving (E47.6).
The geometric kernels in (E47.8) follow from
the actual quotient $F\to F/[I_y,I_y]$
and its restriction to $N$ in E.38; they
retain the arithmetic origin used throughout.

The exact program verifies forty identities
(E47.15) in the integral cyclotomic rings
of conductors $3,9,27,81$, all three norm
identities (E47.19) on their conjugates,
and the finite centered isomorphisms at five
positional levels. At level $n$ the finite
version of $7S_7-1$ has cyclic cokernel
$\mathbb Z_3/(7^{3^{n-1}}-1)$, of order
$3^n$; its evaluation map and every
transition through $n=5$ are checked.
These are exact algebraic checks of the
maps used above. The all-level class-number
and generator-index assertions follow from
(E47.13)--(E47.18).
\end{proof}
